\documentclass{article}
\usepackage[utf8]{inputenc}
\usepackage{geometry}
\geometry{a4paper, margin=1in}
\usepackage{hyperref}
\usepackage{listings}
\usepackage{xcolor}

\definecolor{codegreen}{rgb}{0,0.6,0}
\definecolor{codegray}{rgb}{0.5,0.5,0.5}
\definecolor{codepurple}{rgb}{0.58,0,0.82}
\definecolor{backcolour}{rgb}{0.95,0.95,0.92}

\lstdefinestyle{mystyle}{
    backgroundcolor=\color{backcolour},   
    commentstyle=\color{codegreen},
    keywordstyle=\color{magenta},
    numberstyle=\tiny\color{codegray},
    stringstyle=\color{codepurple},
    basicstyle=\ttfamily\footnotesize,
    breakatwhitespace=false,         
    breaklines=true,                 
    captionpos=b,                    
    keepspaces=true,                 
    numbers=left,                    
    numbersep=5pt,                  
    showspaces=false,                
    showstringspaces=false,
    showtabs=false,                  
    tabsize=2
}
\lstset{style=mystyle}

\title{Technical Execution Plan: Phases 1 \& 2\\
\large Underwater 3D Gaussian Splatting Reconstruction}
\author{Joshua Hecke}
\date{\today}

\begin{document}

\maketitle

\section{Introduction}
This document outlines the detailed technical steps required to execute Phase 1 (Foundation and Method Selection) and Phase 2 (Baseline Reconstruction and Optimization) of the Scene Change Detection in Underwater 3D Gaussian Splatting project. It specifically addresses two critical constraints:
\begin{enumerate}
    \item Utilizing a dual-boot Ubuntu setup on a high-performance machine equipped with an NVIDIA RTX 5070 Ti (16GB GDDR7).
    \item Leveraging the open-sourced \textbf{Submerged3D} dataset for initial testing and validation before accessing the AIMS EcoRRAP dataset.
\end{enumerate}

\section{Phase 1: Environment and Hardware Setup}

\subsection{Ubuntu and NVIDIA Drivers}
The machine utilizes an RTX 5070 Ti, which is based on the latest architecture. It is critical to use the latest NVIDIA proprietary drivers.
\begin{lstlisting}[language=bash]
# After installing Ubuntu 24.04 (or 22.04), update the system
sudo apt update && sudo apt upgrade -y

# Install standard build essentials
sudo apt install build-essential git wget curl -y

# Add the proprietary graphics drivers PPA and install the latest driver
sudo add-apt-repository ppa:graphics-drivers/ppa
sudo apt update
sudo ubuntu-drivers autoinstall
# Reboot the machine after installation
sudo reboot
\end{lstlisting}

\subsection{CUDA Toolkit and Conda Configuration}
Because the RTX 5070 Ti requires up-to-date software, we will aim to use the latest CUDA versions (e.g., CUDA 12.1+). 

\textbf{Important Note on CUDA Compatibility:} Several of the evaluated codebases originally specify \textbf{CUDA 11.8}. If you encounter compilation errors (specifically with the \texttt{diff-gaussian-rasterization} modules) when using CUDA 12.x, you will need to downgrade the environment's PyTorch and CUDA toolkit to 11.8. Notes are provided in each section below.

\begin{lstlisting}[language=bash]
# Download and install Miniconda
mkdir -p ~/miniconda3
wget https://repo.anaconda.com/miniconda/Miniconda3-latest-Linux-x86_64.sh -O ~/miniconda3/miniconda.sh
bash ~/miniconda3/miniconda.sh -b -u -p ~/miniconda3
rm -rf ~/miniconda3/miniconda.sh
~/miniconda3/bin/conda init bash
# Restart terminal
\end{lstlisting}

\subsection{Dataset Preparation (Submerged3D)}
We will use the Submerged3D dataset hosted on HuggingFace.
\begin{lstlisting}[language=bash]
# Install HuggingFace CLI
pip install -U "huggingface_hub[cli]"

# Download the Submerged3D dataset
huggingface-cli download theflash987/Submerged3D --repo-type dataset --local-dir ~/UMP/Codebase/Dataset/Submerged3D
\end{lstlisting}
The Submerged3D dataset contains 4 environments (e.g., Cormoran, Isro, Kwaj, Tokai), each with 20 images in an \texttt{images} folder, but no COLMAP data.

\subsection{Generating COLMAP Data}
The \texttt{convert.py} script from the standard 3DGS repository expects input images to be located in an \texttt{input/} folder. Because the Submerged3D dataset provides images in an \texttt{images/} folder, you must first rename this folder to \texttt{input/}.

\begin{lstlisting}[language=bash]
# Rename the images directory for each scene
mv ~/UMP/Codebase/Dataset/Submerged3D/<scene_name>/images ~/UMP/Codebase/Dataset/Submerged3D/<scene_name>/input
\end{lstlisting}

Then, you can generate the required sparse point cloud and camera poses for 3DGS using the \texttt{convert.py} script:

\begin{lstlisting}[language=bash]
# Install COLMAP (and FFmpeg if needed)
sudo apt install colmap ffmpeg -y

# The original 3D Gaussian Splatting repository is already provided in the Tools directory.
cd ~/UMP/Codebase/Tools/gaussian-splatting-main

# Create a temporary conda environment to run the script
conda create -n colmap_runner python=3.9 -y
conda activate colmap_runner

# Run the convert script on each of the 4 scenes
# Replace <scene_name> with Cormoran, Isro, Kwaj, and Tokai
python convert.py -s ~/UMP/Codebase/Dataset/Submerged3D/<scene_name>
\end{lstlisting}

This will create a \texttt{sparse/0} folder inside each scene directory, containing the \texttt{cameras.bin}, \texttt{images.bin}, and \texttt{points3D.bin} files required by the 3DGS approaches. The script will also create an \texttt{images} folder containing undistorted versions of the inputs.

\subsection{Depth Map Generation and Frame Interpolation}
Several models (3D-UIR, UW-GS, and RUSplatting) strictly require or highly recommend depth maps. RUSplatting additionally requires inverted depth maps and heavily relies on interpolated frames due to the sparse views (only 20 images per scene).

\textbf{1. Depth Map Generation (Depth-Anything-V2)}
\begin{lstlisting}[language=bash]
cd ~/UMP/Codebase/Tools/Depth-Anything-V2-main
conda create -n depth_anything python=3.10 -y
conda activate depth_anything
pip install torch torchvision torchaudio
pip install -r requirements.txt

# Download the weights manually
mkdir checkpoints
wget "https://huggingface.co/depth-anything/Depth-Anything-V2-Large/resolve/main/depth_anything_v2_vitl.pth?download=true" -O checkpoints/depth_anything_v2_vitl.pth

# For standard depth maps (e.g., for 3D-UIR and UW-GS):
python run.py --encoder vitl --pred-only --grayscale --img-path ~/UMP/Codebase/Dataset/Submerged3D/<scene_name>/images --outdir ~/UMP/Codebase/Dataset/Submerged3D/<scene_name>/depthmap

# For inverted depth maps (Required by RUSplatting):
# Note: You MUST modify run.py line 63 to include `depth = 255.0 - depth` before running this.
python run.py --encoder vitl --pred-only --grayscale --img-path ~/UMP/Codebase/Dataset/Submerged3D/<scene_name>/images --outdir ~/UMP/Codebase/Dataset/Submerged3D/<scene_name>/depthmap_inverted
\end{lstlisting}

\textbf{2. Frame Interpolation (RIFE) - Required for RUSplatting}
RUSplatting requires generating intermediate frames between your 20 source images (named with a \texttt{\_to\_} flag, e.g., \texttt{frame000\_to\_frame001.jpg}). You will need to utilize a tool like RIFE (Real-Time Intermediate Flow Estimation) to generate these prior to running RUSplatting.\section{Phase 2: Baseline Reconstruction and Optimization}
We will train 6 different models. For each model, assume you are running the commands from the root of the respective repository inside \texttt{~/UMP/Codebase/3DGS-Water-Approaches/}.

\textbf{Note on Submodules:} Before installing dependencies for any model, run \texttt{git submodule update --init --recursive} inside the repository to ensure all submodules are properly downloaded.

\subsection{Model 1: SeaSplat}
\begin{lstlisting}[language=bash]
cd ~/UMP/Codebase/3DGS-Water-Approaches/Physics/seasplat-master
conda create --name seasplat_py310 -y python=3.10
conda activate seasplat_py310

# Try latest CUDA first. If it fails, use: pip install torch==2.1.0 torchvision==0.16.0 torchaudio==2.1.0 --index-url https://download.pytorch.org/whl/cu118
pip install torch torchvision torchaudio
pip install -r requirements.txt
pip install submodules/diff-gaussian-rasterization
pip install submodules/simple-knn

# Train
python train.py -s ~/UMP/Codebase/Dataset/Submerged3D/<scene_name> --exp seasplat_exp --do_seathru --seathru_from_iter 10000 --eval
\end{lstlisting}

\subsection{Model 2: 3D-UIR}
\begin{lstlisting}[language=bash]
cd ~/UMP/Codebase/3DGS-Water-Approaches/Physics/3D-UIR-main
conda create -n 3d-uir python=3.10
conda activate 3d-uir

# NOTE: Originally uses CUDA 11.8. Try latest first. If failure:
# pip install torch==2.1.0 torchvision==0.16.0 torchaudio==2.1.0 --index-url https://download.pytorch.org/whl/cu118
# conda install cudatoolkit-dev=11.8 -c conda-forge
pip install torch torchvision torchaudio

pip install submodules/diff-gaussian-rasterization
pip install submodules/simple-knn/
pip install submodules/fused-ssim/
pip install ninja git+https://github.com/NVlabs/tiny-cuda-nn/#subdirectory=bindings/torch

# 3D-UIR expects the depth directory to be named "depths" instead of "depthmap"
# It also requires a depth parameters JSON file
cd ~/UMP/Codebase/Dataset/Submerged3D/<scene_name>
ln -s depthmap depths
cd ~/UMP/Codebase/3DGS-Water-Approaches/Physics/3D-UIR-main

# Generate depth parameters for 3D-UIR
python utils/make_depth_scale.py --base_dir ~/UMP/Codebase/Dataset/Submerged3D/<scene_name> --depths_dir ~/UMP/Codebase/Dataset/Submerged3D/<scene_name>/depths

# Train
python train.py -s ~/UMP/Codebase/Dataset/Submerged3D/<scene_name> -d ~/UMP/Codebase/Dataset/Submerged3D/<scene_name>/depths --eval
\end{lstlisting}

\subsection{Model 3: Gaussian Splashing}
\begin{lstlisting}[language=bash]
cd ~/UMP/Codebase/3DGS-Water-Approaches/Image/gaussianSplashing-main
conda env create --file environment.yml
conda activate gaussianSplashing_env

# Train using Hybrid mode
python train.py -s ~/UMP/Codebase/Dataset/Submerged3D/<scene_name> --underwater_processing HYB --eval
\end{lstlisting}

\subsection{Model 4: WaterSplatting}
\begin{lstlisting}[language=bash]
cd ~/UMP/Codebase/3DGS-Water-Approaches/Image/water-splatting-main
conda create --name water_splatting -y python=3.8
conda activate water_splatting

# NOTE: Originally built for CUDA 11.8. 
pip install torch==2.1.2+cu118 torchvision==0.16.2+cu118 --extra-index-url https://download.pytorch.org/whl/cu118
conda install -c "nvidia/label/cuda-11.8.0" cuda-toolkit
pip install ninja git+https://github.com/NVlabs/tiny-cuda-nn/#subdirectory=bindings/torch
pip install nerfstudio==1.1.4
ns-install-cli
pip install --no-use-pep517 -e .

# Train
ns-train water-splatting --vis viewer+wandb colmap --downscale-factor 1 --colmap-path sparse/0 --data ~/UMP/Codebase/Dataset/Submerged3D/<scene_name> --images-path images
\end{lstlisting}

\subsection{Model 5: RUSplatting}
\begin{lstlisting}[language=bash]
cd ~/UMP/Codebase/3DGS-Water-Approaches/Additional/RUSplatting-main
conda create -n rusplatting python=3.12
conda activate rusplatting

# Uses PyTorch 2.5.1 and CUDA 12.8 natively
pip install torch torchvision torchaudio
pip install submodules/diff-gaussian-rasterization
pip install submodules/simple-knn

# RUSplatting expects the inverted depth maps to be available in "depthmap"
cd ~/UMP/Codebase/Dataset/Submerged3D/<scene_name>
mv depthmap depthmap_standard
ln -s depthmap_inverted depthmap
cd ~/UMP/Codebase/3DGS-Water-Approaches/Additional/RUSplatting-main

# Train
python train.py -s ~/UMP/Codebase/Dataset/Submerged3D/<scene_name> --adaptive --eval

# Restore original depthmap after training
cd ~/UMP/Codebase/Dataset/Submerged3D/<scene_name>
rm depthmap
mv depthmap_standard depthmap
\end{lstlisting}

\subsection{Model 6: UW-GS}
\begin{lstlisting}[language=bash]
cd ~/UMP/Codebase/3DGS-Water-Approaches/Additional/UW-GS-main
conda env create -f environment.yml
conda activate UW-GS

# NOTE: Environment targets Python 3.7, PyTorch 1.12.1, CUDA 11.6. It is highly recommended to follow these versions strictly for this codebase.
pip install imageio==2.27.0 opencv-python imageio-ffmpeg scipy dearpygui lpips

# Train
python train.py -s ~/UMP/Codebase/Dataset/Submerged3D/<scene_name> -m output/scene --BMM_Flag --eval
\end{lstlisting}

\section{Mesh Extraction with SuGaR (The 7th Splat)}
Once the 6 models are trained, evaluate their output point clouds visually. Select the output directory of the model that provides the cleanest water-removed geometry. We will use this pre-trained 3DGS model as a prior for SuGaR.

\begin{lstlisting}[language=bash]
cd ~/UMP/Codebase/3DGS-Water-Approaches/Additional/SuGaR-main
conda create --name sugar -y python=3.9
conda activate sugar

# NOTE: Requires CUDA 11.8 for compatibility with PyTorch3D and its specific rasterizer versions.
conda install pytorch==2.0.1 torchvision==0.15.2 torchaudio==2.0.2 pytorch-cuda=11.8 -c pytorch -c nvidia
conda install -c fvcore -c iopath -c conda-forge fvcore iopath
conda install pytorch3d==0.7.4 -c pytorch3d
conda install -c plotly plotly
conda install -c conda-forge rich plyfile==0.8.1 jupyterlab nodejs ipywidgets
pip install open3d PyMCubes

pip install -e gaussian_splatting/submodules/diff-gaussian-rasterization/
pip install -e gaussian_splatting/submodules/simple-knn/

# Run SuGaR on the best pre-trained model
python train_full_pipeline.py -s ~/UMP/Codebase/Dataset/Submerged3D/<scene_name> -r "dn_consistency" --high_poly True --export_obj True --gs_output_dir <path_to_best_model_output>
\end{lstlisting}

This will generate a highly detailed, surface-aligned 3D Gaussian Splatting model and extract a textured \texttt{.obj} mesh representing the clear-water baseline structure.

\section{Evaluation and Quantitative Metrics}
To produce a quantitative table comparing the models (e.g., PSNR, SSIM, LPIPS), it is critical that the models are trained with the \texttt{--eval} flag (as specified in the commands above). This flag holds out a subset of images (usually every 8th image) for testing.

\subsection{Standard 3DGS-based Models}
For SeaSplat, 3D-UIR, Gaussian Splashing, RUSplatting, and UW-GS, evaluation is a two-step process using the scripts provided in the respective forks:

\begin{lstlisting}[language=bash]
# Step 1: Render the test set
# Run this from the root of the specific model's repository
python render.py -m <path_to_model_output> --skip_train

# Step 2: Compute metrics
python metrics.py -m <path_to_model_output>
\end{lstlisting}
The output will be saved in \texttt{<path\_to\_model\_output>/results.json}, detailing the PSNR, SSIM, and LPIPS scores on the holdout set.

\subsection{WaterSplatting (Nerfstudio)}
WaterSplatting automatically computes metrics during training and evaluates on a holdout set by default. You can view the final metrics in the console output or the WandB dashboard (if configured). Alternatively, you can explicitly evaluate a trained model by running:
\begin{lstlisting}[language=bash]
ns-eval --load-config <path_to_outputs>/config.yml
\end{lstlisting}

\subsection{Viewing Outputs Visually}
To visually inspect the trained models:
\begin{itemize}
    \item \textbf{Standard 3DGS Models}: Use the pre-compiled \textbf{SIBR Viewers} provided by the original 3D Gaussian Splatting repository. Alternatively, you can drag and drop the output \texttt{point\_cloud.ply} (located in \texttt{<path\_to\_model\_output>/point\_cloud/iteration\_30000/}) into a WebGL viewer like \href{https://playcanvas.com/supersplat/editor}{SuperSplat} or \href{https://antimatter15.com/splat/}{Antimatter15's Splat Viewer}.
    \item \textbf{WaterSplatting (Nerfstudio)}: Run \texttt{ns-viewer --load-config <path\_to\_outputs>/config.yml} and open the provided localhost link in your browser.
\end{itemize}

\end{document}
